%% tutorial.tex -- Open-BioAFE Design Tutorial (BRICD)
%% Compile with pdfLaTeX + BibTeX (Overleaf: set this file as the Main document).
%% All sources live in the tutorial/ folder.
\documentclass[11pt,a4paper,oneside,openany]{book}
\usepackage{tutorial/biotut}

% Faster compiles while editing: uncomment and list the chapters you are working on;
% cross-references to other chapters keep working.
%\includeonly{tutorial/chapters/02_signals,tutorial/chapters/03_electrodes}
% Partial builds (local): if <jobname>.only exists it may contain \includeonly{...}
\IfFileExists{\jobname.only}{\input{\jobname.only}}{}

\hypersetup{pdftitle={Designing a Biopotential Acquisition Chip from Scratch: The Open-BioAFE Tutorial},
            pdfauthor={Open-BioAFE Project, BRICD}}

\begin{document}
\frontmatter
%% ------------------------------------------------------------------ title page
\begin{titlepage}
\thispagestyle{empty}
\tikzexternaldisable
\begin{tikzpicture}[remember picture,overlay]
  \fill[navy] (current page.south west) rectangle (current page.north east);
  \foreach \x in {0,1,...,21}{\draw[slate,line width=0.3pt] ($(current page.south west)+(\x*1cm,0)$) -- ++(0,30cm);}
  \foreach \y in {0,1,...,30}{\draw[slate,line width=0.3pt] ($(current page.south west)+(0,\y*1cm)$) -- ++(21cm,0);}
  % stylised ECG trace: three beats (P wave, QRS complex, T wave)
  \begin{scope}[shift={($(current page.south west)+(1.2cm,6.2cm)$)},x=1cm,y=1cm]
    \foreach \s in {0,6.2,12.4}{
      \draw[goldlt,line width=2.2pt,line cap=round,line join=round]
        (\s,0) -- (\s+0.9,0) .. controls (\s+1.15,0.55) and (\s+1.45,0.55) .. (\s+1.7,0)
        -- (\s+2.15,0) -- (\s+2.35,-0.45) -- (\s+2.65,3.6) -- (\s+2.95,-0.9) -- (\s+3.15,0)
        -- (\s+3.7,0) .. controls (\s+4.1,1.05) and (\s+4.7,1.05) .. (\s+5.1,0) -- (\s+6.2,0);
    }
  \end{scope}
  \node[anchor=north west,text=goldlt,font=\sffamily\bfseries\large] at ($(current page.north west)+(2.4cm,-3.2cm)$)
     {OPEN-BIOAFE \textperiodcentered{} IHP SG13G2 \textperiodcentered{} 8-CHANNEL ExG FRONT-END};
  \node[anchor=north west,text=white,font=\sffamily\bfseries\fontsize{40}{46}\selectfont,text width=16.5cm,align=left]
     at ($(current page.north west)+(2.4cm,-4.4cm)$) {Designing a Biopotential Acquisition Chip from Scratch};
  \draw[gold,line width=2.4pt] ($(current page.north west)+(2.45cm,-10.3cm)$) -- ++(4.5cm,0);
  \node[anchor=north west,text=pale,font=\sffamily\Large,text width=15.5cm,align=left]
     at ($(current page.north west)+(2.4cm,-10.9cm)$)
     {A design tutorial for undergraduate and early-graduate engineers building an open-source biomedical analog front-end and data converters in 130\,nm SiGe BiCMOS};
  \node[anchor=north west,text=white,font=\sffamily\normalsize,text width=15.5cm,align=left]
     at ($(current page.north west)+(2.4cm,-14.6cm)$)
     {Bioelectric signals \textperiodcentered{} Electrodes and safety \textperiodcentered{} Low-noise amplifiers \textperiodcentered{} Chopping \textperiodcentered{} SAR and $\Delta\Sigma$ ADCs \textperiodcentered{} Bioimpedance \textperiodcentered{} Biosignal processing \textperiodcentered{} Open-source flow \textperiodcentered{} Hands-on labs};
  \node[anchor=south west,text=white,font=\sffamily\bfseries\large] at ($(current page.south west)+(2.4cm,2.6cm)$) {BRICD};
  \node[anchor=south west,text=pale,font=\sffamily\small] at ($(current page.south west)+(2.4cm,1.9cm)$)
     {Biomedical Research in Integrated Circuit Design \textperiodcentered{} thebricd.org};
  \node[anchor=south east,text=pale,font=\sffamily\small,align=right] at ($(current page.south east)+(-2.4cm,1.9cm)$)
     {Open-BioAFE Project\\Draft v0.1 \textperiodcentered{} November 2026};
\end{tikzpicture}
\tikzexternalenable
\end{titlepage}

\tikzsetfigurename{00_preface-}
%% 00_preface.tex -- Preface (front matter, unnumbered)
\chapter{Preface}\label{ch:preface}

A heartbeat on the skin is about a millivolt. A brain rhythm on the scalp is a few tens of microvolts. A single neuron, seen from an electrode a few tens of micrometres away, is a hundred microvolts that lasts a millisecond. Every wearable ECG patch, EEG headset, prosthetic-control armband and neural implant has to turn these signals into clean numbers while the patient's body picks up volts of mains interference, the electrodes add hundreds of millivolts of offset, and the battery allows a few microwatts per channel. The chip that does this is a \textbf{biopotential analog front-end} (AFE) and its data converters.

This book teaches you to design one by following a single, concrete chip, \textbf{\OB{}}: an 8-channel ECG/EEG/EMG front-end with a 12-bit SAR and a 16-bit incremental $\Delta\Sigma$ converter, built in the open-source IHP \sgx{} \qty{130}{\nano\meter} SiGe BiCMOS process~\cite{ihppdk} with open-source tools~\cite{iicosic}. Every concept follows the same pattern as our \emph{Open-SERDES} tutorial~\cite{serdestut}: a short derivation, a worked example with the project's real numbers, and a \emph{what this means for your design} box.

\paragraph{Who this book is for} Third- and fourth-year electrical engineering undergraduates and first-year graduate students, and biomedical engineers who need to specify the hardware they will use. After reading it you should be able to:
\begin{itemize}
  \item turn a clinical or research application into a front-end specification: channels, bandwidth, noise, CMRR, input impedance, offset range and power;
  \item design and budget a low-noise, low-power amplifier chain: chopping, DC servo loops, filters and right-leg drive;
  \item choose, size and verify a SAR or $\Delta\Sigma$ converter for a biosignal;
  \item show, with real ECG, EEG and EMG recordings, how the circuit choices move a clinical result;
  \item carry a block from schematic to a verified layout with open tools, ready for an IHP OpenMPW shuttle.
\end{itemize}

\paragraph{Prerequisites} Circuit analysis, signals and systems (Fourier and Laplace transforms, filters) and one course in analog electronics (small-signal MOS and BJT models, the differential pair, op-amps), plus some Python. The transistor-level background of Part~II, in particular the $\gm/I_\mathrm{D}$ method, is taught in Part~II of \emph{The Art of Analog IC Design}~\cite{artofaicd}; we cross-link it instead of repeating it.

\section*{How the book is organised}
The book has five parts. \textbf{Part~I, Foundations} (Chapters~1--4), explains where biopotentials come from, what the electrodes and the body do to them, and the interference, safety rules and standards every front-end must meet. \textbf{Part~II, The Analog Front-End} (Chapters~5--12), designs the low-noise, low-power amplifier chain. \textbf{Part~III, Data Converters} (Chapters~13--22), is the heart of the book: it takes you from the quantiser to a measured SAR and $\Delta\Sigma$ converter. \textbf{Part~IV, Biosignal Processing and Systems} (Chapters~23--28), processes real ECG, EEG and EMG data, puts processing on chip and adds bioimpedance sensing. \textbf{Part~V, Putting It Together} (Chapters~29--31), assembles the architecture and its budgets, walks through the open-source flow and collects the fifteen labs.

\paragraph{Three reading paths} Everyone should read Part~I first. After that, \cref{fig:pre-paths} suggests three routes. The \textbf{circuit path} (Parts I, II, III and V) is for students heading to a tape-out. The \textbf{signal path} (Part~I, Part~IV, the converter specifications of \cref{ch:sampling,ch:adcspecs}, then Part~V) is for biomedical-engineering and DSP students who need to specify, not build, the hardware. The \textbf{converter-only path} (\cref{ch:sampling} to \cref{ch:adctest} with Labs~6--10) is a self-contained course on data converters.

\begin{figure}[htbp]
\centering
\tikzset{pnode/.style={draw=#1,fill=white,line width=0.8pt,rounded corners=2pt,minimum height=7mm,minimum width=13mm,
    inner sep=2pt,font=\sffamily\footnotesize\bfseries,text=#1,align=center},
  plink/.style={line width=1.1pt,draw=#1}}
\begin{tikzpicture}[x=1cm,y=1cm]
\node[pnode=navy] (p1) at (0,0) {Part I\\[-2pt]\scriptsize Ch 1--4};
% circuit path
\node[pnode=golddk] (c2) at (2.6,1.3) {Part II\\[-2pt]\scriptsize Ch 5--12};
\node[pnode=golddk] (c3) at (5.0,1.3) {Part III\\[-2pt]\scriptsize Ch 13--22};
% signal path
\node[pnode=bblue] (s4) at (2.6,0) {Part IV\\[-2pt]\scriptsize Ch 23--28};
\node[pnode=bblue] (s13) at (5.0,0) {Ch 13--14\\[-2pt]\scriptsize ADC specs};
% converter-only path
\node[pnode=slate] (k1) at (2.6,-1.3) {Ch 13--22};
\node[pnode=slate] (k2) at (5.0,-1.3) {Labs 6--10};
% merge
\node[pnode=navy] (p5) at (7.6,0.65) {Part V\\[-2pt]\scriptsize Ch 29--31};
\draw[plink=golddk] (p1) -- (c2); \draw[plink=golddk] (c2) -- (c3); \draw[plink=golddk] (c3) -- (p5);
\draw[plink=bblue] (p1) -- (s4); \draw[plink=bblue] (s4) -- (s13); \draw[plink=bblue] (s13) -- (p5);
\draw[plink=slate] (k1) -- (k2);
\node[font=\sffamily\scriptsize\bfseries,text=golddk,anchor=south] at (3.8,1.7) {circuit path: to a tape-out};
\node[font=\sffamily\scriptsize\bfseries,text=bblue,anchor=south] at (3.8,0.36) {signal path: specify the hardware};
\node[font=\sffamily\scriptsize\bfseries,text=slate,anchor=north] at (3.8,-1.72) {converter-only path: a data-converter course};
\end{tikzpicture}
\caption{\textbf{Three ways to read the book.} The circuit path (gold) follows the front-end and converter circuits to a tape-out. The signal path (blue) goes from the signals to their processing and the converter specifications. The converter-only path (grey) is Part~III with its labs on its own.}
\label{fig:pre-paths}
\end{figure}

\section*{How to read the boxes}
Chapters open with what you will learn and close with key takeaways. In between, six kinds of boxes mark the role of each piece of text:

{\renewcommand{\thewexample}{P.1}%
\begin{keyeq}[Signal-to-noise ratio in decibels]
\[ \SNR = 20\log_{10}\frac{V_\mathrm{signal,rms}}{V_\mathrm{noise,rms}} \]
\end{keyeq}
\begin{example}{A number you will use on every page}
A \qty{20}{\micro\volt} rms EEG rhythm over the \qty{1.5}{\uVrms} noise target of \OB{} gives $\SNR = 20\log_{10}(20/1.5) \approx \qty{22.5}{\deci\bel}$.
\end{example}}
\setcounter{wexample}{0}
\begin{insight}
An \emph{insight} box says what a result means for the \OB{} design.
\end{insight}
\begin{pitfall}
A \emph{pitfall} box warns about a common mistake, such as quoting an amplitude in rms when the specification is peak-to-peak.
\end{pitfall}
\begin{tryit}[Lab 1: biosignal explorer]
A \emph{try-it} box points to a lab script (\cref{sec:lab1}).
\end{tryit}
\begin{specnote}
A \emph{specnote} box states an \OB{} target: input-referred noise below \qty{1.5}{\uVrms} from \qtyrange{0.5}{100}{\hertz} in EEG mode.
\end{specnote}
Longer algebra sits in grey \emph{derivation} boxes.

\section*{How the labs work}
Fifteen labs in three tracks use only free tools: Python (NumPy, SciPy, \tool{wfdb}, \tool{python-deltasigma}) for signals and system models, \tool{Xschem} and \tool{ngspice} for transistor-level \sgx{} circuits, and \tool{KLayout} for layout (\cref{ch:labs}). The Python labs live in \code{labs/python/} and share one module, \code{biosig.py}, which holds the signal tables, the synthetic signals and the electrode and interference models. The figure scripts in \code{tutorial/figsrc/} import the same module, so the numbers in the book and the numbers your labs print come from the same code. Datasets stay on PhysioNet~\cite{physionet} and NinaPro~\cite{ninapro} under their own licences; the labs download them on first use.

\section*{The state of this draft}
The book is being written in six stages. This first release (v0.1) contains the complete outline of all 31 chapters, the four chapters of Part~I and Labs~1 and~2 in full, and the appendices. Chapters still in outline form say so at the top and list their topics, their design question and their labs. The next stage adds the Python-only converter core (\cref{ch:sampling,ch:adcspecs,ch:dstheory} with Labs~6 and~9).

\section*{Sources and acknowledgments}
The biomedical-instrumentation chapters draw on Webster's \emph{Medical Instrumentation}~\cite{webster} and on Yazicioglu, Van Hoof and Puers' \emph{Biopotential Readout Circuits}~\cite{yazicioglubook}; the converter chapters on Pelgrom~\cite{pelgrom} and on Schreier, Temes and Pavan~\cite{schreier}. The chip is possible because IHP opened its \sgx{} process design kit~\cite{ihppdk} and the Institute for Integrated Circuits at JKU Linz maintains the IIC-OSIC-TOOLS container and the analog chip template~\cite{iicosic,jkutemplate}. The labs use open data from PhysioNet~\cite{physionet} and NinaPro~\cite{ninapro}.

This is a \textbf{living draft (v0.1)}; numbers will be refined as the design matures. If you find an error, an unclear derivation or a figure that does not explain what it should, please email \href{mailto:bricd.org@gmail.com}{bricd.org@gmail.com} or open an issue or a pull request on the project's GitHub repository, linked from thebricd.org.

\section*{How this book was made}
This book, its figures and its lab code were developed with the help of AI tools, and this draft may contain mistakes; please report any you find as described above, and we will correct them. The topics, structure and many of the results follow the sources acknowledged above and listed in the References. Standards are summarised, never reproduced: always work from the edition of IEC~60601 that applies to your device.

{\hypersetup{linkcolor=navy}\tableofcontents}

\mainmatter
\tutpart{Foundations}{Where biopotentials come from, how big and how fast they are, what the electrode and the body do to them, and the interference, safety rules and standards every front-end must meet. Every budget in the rest of the tutorial starts here.}
\include{tutorial/chapters/01_why}
\include{tutorial/chapters/02_signals}
\include{tutorial/chapters/03_electrodes}
\include{tutorial/chapters/04_interference}

\tutpart{The Analog Front-End}{Low-noise, low-power amplification at the microwatt level: devices in weak inversion, the noise--power trade-off, instrumentation amplifiers, the AC-coupled LNA, chopping, DC servo loops, filters and common-mode control, all in SG13G2.}
%% 05_devices.tex -- Chapter 5: Devices for Microwatt Design
%% Status: OUTLINE (build stage 3). Replace the outline with the full chapter when it is written.
\chapter{Devices for Microwatt Design}\label{ch:devices}
\chapterlead{A biopotential amplifier lives on a few hundred nanoamperes per transistor, so its devices sit in weak or moderate inversion, where the square law no longer applies.}

\begin{specnote}[Chapter status: outline]
This chapter is an outline. It will be written in build stage 3 (Part II with Labs 3--5). The topics, the design question and the labs below are fixed; the derivations, worked examples and figures are still to come.
\end{specnote}

\begin{objectives}
\begin{itemize}
  \item place a MOSFET in weak, moderate or strong inversion from its inversion coefficient and read its $\gm/I_\mathrm{D}$;
  \item explain why $\gm/I_\mathrm{D}$ peaks at about $1/(n U_T)$ in weak inversion and what that means for noise per microampere;
  \item relate flicker noise to gate area and pick input-pair dimensions for a sub-\qty{100}{\hertz} band;
  \item compare \sgx{} MOSFETs and SiGe HBTs as low-noise input devices (thermal noise, flicker noise, base current);
  \item estimate offset and gain mismatch from Pelgrom coefficients.
\end{itemize}
\end{objectives}

\begin{insight}[Design question]
Which device, at what inversion level and size, gives the lowest input-referred noise for a microwatt bias current?
\end{insight}

\section{Weak, moderate and strong inversion}\label{sec:ch5-1}
The inversion coefficient $IC = I_\mathrm{D}/I_\mathrm{spec}$ places a transistor on one continuous curve from exponential (weak inversion) to square-law (strong inversion) behaviour. Biomedical front-ends bias their input pairs at $IC \lesssim 1$.

\section{The \texorpdfstring{$\gm/I_\mathrm{D}$}{gm/ID} method at microwatt bias}\label{sec:ch5-2}
Sizing from $\gm/I_\mathrm{D}$ lookup tables instead of hand equations; the weak-inversion limit $\gm/I_\mathrm{D} \approx 1/(nU_T) \approx \qty{25}{\per\volt}$; why it saturates and what it costs in area and speed. This chapter cross-links \emph{The Art of Analog IC Design}~\cite{artofaicd} rather than repeating it.

\section{Flicker noise and device area}\label{sec:ch5-3}
$1/f$ noise scales as $1/(WL)$, so the low corner of an EEG band forces large input devices; PMOS versus NMOS in \sgx{}; the corner frequency and its dependence on bias.

\section{MOSFET versus SiGe HBT input pairs}\label{sec:ch5-4}
An HBT reaches $\gm = I_\mathrm{C}/U_T$ at any current and has very low $1/f$ noise, but its base current flows through the electrode impedance and adds shot noise. The trade-off decides the input stage of the \OB{} LNA.

\section{Matching}\label{sec:ch5-5}
Pelgrom's law $\sigma(\Delta V_\mathrm{T}) = A_{V_T}/\sqrt{WL}$, current-mirror and capacitor matching, and why large, slow biomedical devices match well but still need chopping for microvolt offsets.

\begin{tryit}[Labs for this chapter]
Lab~3 (\cref{sec:lab3}) characterises $\gm/I_\mathrm{D}$, flicker noise versus $W\!\cdot\!L$ and HBT versus MOS input noise in \sgx{}; its lookup tables are reused in Labs~4 and~5.
\end{tryit}

\paragraph{Read with this chapter} Sarpeshkar~\cite{sarpeshkar} (subthreshold design, noise--power limits), Razavi~\cite{razavi} (device noise), \emph{The Art of Analog IC Design}~\cite{artofaicd} (Part II: the transistor and $\gm/I_\mathrm{D}$).

%% 06_noise.tex -- Chapter 6: Noise and the Noise--Power Trade-off
%% Status: OUTLINE (build stage 3). Replace the outline with the full chapter when it is written.
\chapter{Noise and the Noise--Power Trade-off}\label{ch:noise}
\chapterlead{Every microvolt of noise you remove costs current. The noise efficiency factor tells you how close a design is to the physical limit.}

\begin{specnote}[Chapter status: outline]
This chapter is an outline. It will be written in build stage 3 (Part II with Labs 3--5). The topics, the design question and the labs below are fixed; the derivations, worked examples and figures are still to come.
\end{specnote}

\begin{objectives}
\begin{itemize}
  \item compute thermal, flicker, shot and $kT/C$ noise and refer them to the amplifier input;
  \item integrate a noise spectrum over the signal band with the right noise bandwidth;
  \item derive the noise efficiency factor (NEF) and the power efficiency factor (PEF) and use them to compare designs;
  \item explain why a single bipolar or subthreshold device has $\NEF \approx 1$ and a differential pair about~2;
  \item set a current budget per channel from the \qty{1.5}{\uVrms} noise target.
\end{itemize}
\end{objectives}

\begin{insight}[Design question]
What is the minimum supply current that can meet the EEG-mode noise target, and how close to it can a real amplifier get?
\end{insight}

\section{Noise sources}\label{sec:ch6-1}
Thermal noise of resistors and channels, $1/f$ noise, shot noise of junction currents and $kT/C$ noise of sampled capacitors, with their spectral densities and the conditions where each dominates in a biopotential amplifier.

\begin{keyeq}[Noise efficiency factor]
\begin{equation}
  \NEF = V_\mathrm{ni,rms}\sqrt{\frac{2\,I_\mathrm{tot}}{\pi\, U_T\, 4kT\, \mathrm{BW}}}, \qquad \PEF = \NEF^2\,\Vdd
  \label{eq:nef}
\end{equation}
\end{keyeq}

\section{Input-referred noise and noise bandwidth}\label{sec:ch6-2}
Referring output noise to the input; the equivalent noise bandwidth of first- and second-order roll-offs ($\pi/2$ and $\pi/(2\sqrt2)$ factors); integrating $1/f$ noise between a sub-hertz high-pass and the low-pass corner.

\section{Noise efficiency factor and power efficiency factor}\label{sec:ch6-3}
$\NEF = V_\mathrm{ni,rms}\sqrt{2 I_\mathrm{tot}/(\pi U_T\, 4kT\, \mathrm{BW})}$, introduced by Steyaert and Sansen~\cite{steyaert1987} and made the standard yardstick by Harrison and Charles~\cite{harrison2003}; $\PEF = \NEF^2\Vdd$ adds the supply voltage. Published amplifiers plotted against both.

\section{The physical limit}\label{sec:ch6-4}
Why a single BJT or weak-inversion MOSFET gives $\NEF = 1$, a differential pair $\sqrt2$ to~2, and current-reuse topologies below~2. What the limit means for the \qty{10}{\micro\watt} per channel budget.

\begin{tryit}[Labs for this chapter]
Lab~4 (\cref{sec:lab4}) computes the NEF of your Harrison LNA versus bias current.
\end{tryit}

\paragraph{Read with this chapter} Steyaert and Sansen~\cite{steyaert1987}, Harrison and Charles~\cite{harrison2003}, Sarpeshkar~\cite{sarpeshkar}, Razavi~\cite{razavi} (Chapter on noise).

%% 07_ia.tex -- Chapter 7: Instrumentation Amplifier Architectures
%% Status: OUTLINE (build stage 3). Replace the outline with the full chapter when it is written.
\chapter{Instrumentation Amplifier Architectures}\label{ch:ia}
\chapterlead{The instrumentation amplifier is where the common-mode rejection of the whole chain is won or lost.}

\begin{specnote}[Chapter status: outline]
This chapter is an outline. It will be written in build stage 3 (Part II with Labs 3--5). The topics, the design question and the labs below are fixed; the derivations, worked examples and figures are still to come.
\end{specnote}

\begin{objectives}
\begin{itemize}
  \item analyse the three-op-amp, current-feedback and capacitively coupled instrumentation amplifiers (CCIA);
  \item trace where CMRR comes from in each: resistor or capacitor mismatch, finite common-mode gain, input impedance;
  \item design for a rail-to-rail common-mode range at a \qty{1.2}{\volt} supply;
  \item choose an IA architecture for the \OB{} channel.
\end{itemize}
\end{objectives}

\begin{insight}[Design question]
Which instrumentation-amplifier architecture gives more than \qty{90}{\deci\bel} of CMRR and a \qty{\pm 300}{\milli\volt} offset range at microwatt power and \qty{1.2}{\volt}?
\end{insight}

\section{The three-op-amp instrumentation amplifier}\label{sec:ch7-1}
Gain in the first stage, CMRR set by resistor matching in the difference amplifier, and why the classic textbook IA is too power-hungry and offset-limited for a microwatt ExG chip.

\section{Current-feedback (indirect current feedback) IAs}\label{sec:ch7-2}
Transconductors convert the input and feedback voltages to currents; CMRR no longer depends on resistor matching; limited input range and linearity.

\section{Capacitively coupled IAs}\label{sec:ch7-3}
Input capacitors block electrode DC and set the gain by a capacitor ratio, which matches well on chip. CCIAs dominate modern ExG front-ends and lead to the chopper amplifiers of \cref{ch:chopper}.

\section{CMRR mechanisms and mismatch}\label{sec:ch7-4}
Intrinsic CMRR of the amplifier versus CMRR lost to electrode impedance imbalance (\cref{ch:interference}); budgeting the two against the \qty{89}{\deci\bel} the ECG standards imply.

\section{Common-mode range at \texorpdfstring{\qty{1.2}{\volt}}{1.2 V}}\label{sec:ch7-5}
Input stages that tolerate the body common-mode voltage and the \qty{\pm 300}{\milli\volt} electrode offset on a low supply.

\paragraph{Read with this chapter} Yazicioglu, Van Hoof and Puers~\cite{yazicioglubook}, Pall\`as-Areny and Webster~\cite{pallas1991}, Webster~\cite{webster} (biopotential amplifiers).

%% 08_lna.tex -- Chapter 8: The AC-Coupled Biopotential LNA
%% Status: OUTLINE (build stage 3). Replace the outline with the full chapter when it is written.
\chapter{The AC-Coupled Biopotential LNA}\label{ch:lna}
\chapterlead{Two capacitors, two pseudo-resistors and an OTA: the Harrison amplifier is the circuit every neural and ExG front-end is compared with.}

\begin{specnote}[Chapter status: outline]
This chapter is an outline. It will be written in build stage 3 (Part II with Labs 3--5). The topics, the design question and the labs below are fixed; the derivations, worked examples and figures are still to come.
\end{specnote}

\begin{objectives}
\begin{itemize}
  \item derive the gain, band edges and noise of the capacitive-feedback LNA;
  \item explain how MOS-bipolar pseudo-resistors reach high-pass corners in the millihertz range;
  \item design the OTA (folded cascode or current-scaled telescopic) for a given noise and bandwidth;
  \item budget gain, bandwidth and noise for ECG, EEG and EMG modes;
  \item recognise the linearity, leakage and PVT problems of pseudo-resistors.
\end{itemize}
\end{objectives}

\begin{insight}[Design question]
How do you size the capacitors, pseudo-resistors and OTA so that one AC-coupled LNA meets the noise, bandwidth and power targets of every \OB{} mode?
\end{insight}

\section{The Harrison topology}\label{sec:ch8-1}
Midband gain $A_\mathrm{M} = C_1/C_2$, low corner $f_\mathrm{L} = 1/(2\pi R_\mathrm{pr} C_2)$ and high corner set by the OTA transconductance and load; the original design reaches \qtyrange{0.025}{7200}{\hertz} with \qty{2.2}{\uVrms} of noise at \qty{80}{\micro\watt}~\cite{harrison2003}.

\section{Pseudo-resistors}\label{sec:ch8-2}
Diode-connected MOS-bipolar devices give teraohm small-signal resistance in a tiny area; their resistance depends strongly on voltage, temperature and leakage, so the high-pass corner moves.

\section{OTA design}\label{sec:ch8-3}
Folded cascode versus current-scaled telescopic OTAs; putting most of the current into the input pair to approach the NEF limit of \cref{ch:noise}; capacitive noise gain $(C_1+C_2+C_\mathrm{in})/C_2$.

\section{Gain, bandwidth and noise budget}\label{sec:ch8-4}
Allocating the \qty{1.5}{\uVrms} EEG-mode noise target, choosing $C_1$ for the input impedance it presents (\cref{ch:electrodes}), and switching bandwidth between modes.

\section{Linearity and pseudo-resistor problems}\label{sec:ch8-5}
Distortion at large EMG and artefact amplitudes, leakage at high temperature, and tunable alternatives (bias-controlled pseudo-resistors, duty-cycled resistors).

\begin{tryit}[Labs for this chapter]
Lab~4 (\cref{sec:lab4}) designs the capacitive-feedback LNA with pseudo-resistors in \sgx{}: AC, noise and transient simulations and the NEF versus bias current.
\end{tryit}

\paragraph{Read with this chapter} Harrison and Charles~\cite{harrison2003}, Yazicioglu, Van Hoof and Puers~\cite{yazicioglubook}.

%% 09_chopper.tex -- Chapter 9: Chopper Stabilisation and Offset Cancellation
%% Status: OUTLINE (build stage 3). Replace the outline with the full chapter when it is written.
\chapter{Chopper Stabilisation and Offset Cancellation}\label{ch:chopper}
\chapterlead{Chopping moves the signal above the flicker-noise corner, amplifies it there and brings it back, leaving the offset and $1/f$ noise behind.}

\begin{specnote}[Chapter status: outline]
This chapter is an outline. It will be written in build stage 3 (Part II with Labs 3--5). The topics, the design question and the labs below are fixed; the derivations, worked examples and figures are still to come.
\end{specnote}

\begin{objectives}
\begin{itemize}
  \item explain chopping as modulation and demodulation in the frequency domain;
  \item predict the residual offset and the output ripple of a chopper amplifier;
  \item design a ripple-reduction loop;
  \item compare chopping, auto-zeroing and correlated double sampling;
  \item quantify the input-impedance cost of an input chopper.
\end{itemize}
\end{objectives}

\begin{insight}[Design question]
How much offset, ripple and input-impedance loss does chopping leave behind, and how do you suppress each?
\end{insight}

\section{The modulation view of chopping}\label{sec:ch9-1}
Multiplying by a square wave at $f_\mathrm{ch}$ moves the signal to odd harmonics of $f_\mathrm{ch}$; the amplifier's offset and $1/f$ noise stay at baseband and are modulated away by the output chopper.

\section{Residual offset and ripple}\label{sec:ch9-2}
Charge injection and clock feedthrough of the input switches create a residual offset; the up-modulated amplifier offset appears as a triangular ripple at the output.

\section{Ripple-reduction loops}\label{sec:ch9-3}
Sensing the ripple after the output chopper and feeding a correction back; the 60\,nV/$\sqrt{\mathrm{Hz}}$, \qty{60}{\micro\watt} front-end of Yazicioglu et al.~\cite{yazicioglu2007} as the reference design.

\section{Auto-zero and correlated double sampling}\label{sec:ch9-4}
Sampling the offset and subtracting it; the noise folding that makes auto-zero noisier than chopping.

\section{Chopper spikes and input impedance}\label{sec:ch9-5}
An input chopper in front of the input capacitors draws charge every cycle, so the input impedance falls to roughly $1/(2 f_\mathrm{ch} C_\mathrm{in})$; why this matters for dry electrodes and how impedance boosting (\cref{ch:dsl}) recovers it.

\begin{tryit}[Labs for this chapter]
Lab~5 (\cref{sec:lab5}) adds input and output choppers and a DC servo loop to your LNA: $1/f$ removal, ripple, input-impedance loss and \qty{\pm 300}{\milli\volt} offset tolerance.
\end{tryit}

\paragraph{Read with this chapter} Yazicioglu et al.~\cite{yazicioglu2007}, Yazicioglu, Van Hoof and Puers~\cite{yazicioglubook}.

%% 10_dsl.tex -- Chapter 10: Input Impedance, DC Servo Loops and Offset Rejection
%% Status: OUTLINE (build stage 3). Replace the outline with the full chapter when it is written.
\chapter{Input Impedance, DC Servo Loops and Offset Rejection}\label{ch:dsl}
\chapterlead{Electrodes bring hundreds of millivolts of DC and megaohms of source impedance. The front-end must reject the first without loading the second.}

\begin{specnote}[Chapter status: outline]
This chapter is an outline. It will be written in build stage 3 (Part II with Labs 3--5). The topics, the design question and the labs below are fixed; the derivations, worked examples and figures are still to come.
\end{specnote}

\begin{objectives}
\begin{itemize}
  \item derive the input impedance of a capacitively coupled, chopped front-end;
  \item design a positive-feedback impedance-boosting loop;
  \item design a DC servo loop (DSL) that rejects \qty{\pm 300}{\milli\volt} of electrode offset;
  \item compare DC-coupled and AC-coupled front-ends;
  \item state the trade-offs for dry electrodes.
\end{itemize}
\end{objectives}

\begin{insight}[Design question]
How do you tolerate \qty{\pm 300}{\milli\volt} of electrode offset and keep the input impedance high enough for dry electrodes?
\end{insight}

\section{Input impedance of chopped CCIAs}\label{sec:ch10-1}
Switched-capacitor input resistance, its frequency dependence, and the \qty{2.5}{\mega\ohm} the ECG standards require (\cref{ch:interference}).

\section{Positive-feedback impedance boosting}\label{sec:ch10-2}
Pre-charging the input capacitors from the output to cancel the current drawn from the electrode; stability and matching limits.

\section{The DC servo loop}\label{sec:ch10-3}
An integrator in a feedback path subtracts the electrode offset at the input; the loop sets the high-pass corner and must not add noise in band.

\section{DC-coupled versus AC-coupled front-ends}\label{sec:ch10-4}
DC coupling keeps the electrode offset for motion-artefact detection and needs a wide-range ADC; AC coupling removes it in the analog domain. The choice decides between the two converters of Part~III.

\section{Dry electrodes}\label{sec:ch10-5}
Megaohm source impedance, larger offsets and motion artefacts; what each front-end option buys and costs.

\begin{tryit}[Labs for this chapter]
Lab~5 (\cref{sec:lab5}) measures the input impedance with and without a chopper and verifies \qty{\pm 300}{\milli\volt} offset tolerance with the DC servo loop.
\end{tryit}

\paragraph{Read with this chapter} Yazicioglu, Van Hoof and Puers~\cite{yazicioglubook}, Spinelli et al.~\cite{spinelli2003}, Chi et al.~\cite{chi2010}.

%% 11_filters.tex -- Chapter 11: Filters, PGAs and Anti-Aliasing
%% Status: OUTLINE (build stage 3). Replace the outline with the full chapter when it is written.
\chapter{Filters, PGAs and Anti-Aliasing}\label{ch:filters}
\chapterlead{Between the LNA and the ADC the signal needs more gain, a defined bandwidth and protection from aliasing, at corner frequencies from a fraction of a hertz to a few kilohertz.}

\begin{specnote}[Chapter status: outline]
This chapter is an outline. It will be written in build stage 3 (Part II with Labs 3--5). The topics, the design question and the labs below are fixed; the derivations, worked examples and figures are still to come.
\end{specnote}

\begin{objectives}
\begin{itemize}
  \item design Gm-C and switched-capacitor filters with sub-hertz to kilohertz corners;
  \item design a programmable-gain amplifier for the three \OB{} modes;
  \item choose the anti-aliasing filter order from the ADC oversampling ratio;
  \item explain why mains notch filters belong in the digital domain.
\end{itemize}
\end{objectives}

\begin{insight}[Design question]
How much filtering and gain does each mode need before the SAR and the $\Delta\Sigma$ ADC, and where should each filter live?
\end{insight}

\section{Very-low-frequency filters on chip}\label{sec:ch11-1}
Gm-C filters with nanoampere transconductors and switched-capacitor filters at low clock rates; area, noise and linearity at sub-hertz corners.

\section{Programmable gain}\label{sec:ch11-2}
Capacitive and resistive PGAs; gain steps for ECG, EEG and EMG amplitudes (\cref{ch:signals}) so the ADC range is used well.

\section{Anti-aliasing versus oversampling}\label{sec:ch11-3}
The required attenuation at $\fs - f_\mathrm{B}$ for the SAR; why the $\Delta\Sigma$ converter relaxes the analog filter; mains harmonics and EMG folding into the ECG band (\cref{ch:sampling}).

\section{Mains notch filtering}\label{sec:ch11-4}
Analog notches drift and distort; adaptive and comb notches after the ADC (\cref{ch:dsp}) track the mains frequency.

\paragraph{Read with this chapter} Razavi~\cite{razavi} (switched-capacitor circuits), Yazicioglu, Van Hoof and Puers~\cite{yazicioglubook}.

%% 12_cmcontrol.tex -- Chapter 12: Common-Mode Control and Electrode Monitoring
%% Status: OUTLINE (build stage 3). Replace the outline with the full chapter when it is written.
\chapter{Common-Mode Control and Electrode Monitoring}\label{ch:cmcontrol}
\chapterlead{The third electrode is not a reference: it is the output of a feedback loop that pins the body to the amplifier's common mode.}

\begin{specnote}[Chapter status: outline]
This chapter is an outline. It will be written in build stage 3 (Part II with Labs 3--5). The topics, the design question and the labs below are fixed; the derivations, worked examples and figures are still to come.
\end{specnote}

\begin{objectives}
\begin{itemize}
  \item design a right-leg drive (driven-right-leg) loop and analyse its stability;
  \item relate RLD loop gain to the residual common-mode and differential interference;
  \item implement common-mode feedback for the multichannel front-end;
  \item detect lead-off with DC and AC methods;
  \item protect the inputs against ESD and defibrillation pulses in \sgx{}.
\end{itemize}
\end{objectives}

\begin{insight}[Design question]
How much right-leg-drive gain and which protection network let the \OB{} channel pass the ECG CMRR test with margin?
\end{insight}

\section{The right-leg drive loop}\label{sec:ch12-1}
Sensing the average input common mode and driving it back into the body; the loop gain reduces the body common-mode voltage of \cref{ch:interference}; stability with the electrode and body capacitances~\cite{winter1983}.

\section{Common-mode feedback}\label{sec:ch12-2}
Setting the output common mode of the fully differential stages across eight channels.

\section{Lead-off detection}\label{sec:ch12-3}
DC methods (small pull-up/pull-down currents) and AC methods (an excitation above the signal band); leakage-current limits from \cref{ch:interference}; the full impedance measurement follows in \cref{ch:bioz}.

\section{Input protection and ESD}\label{sec:ch12-4}
Series resistors, clamps and their noise and leakage cost; defibrillation-proof inputs; the ESD devices available in \sgx{}.

\begin{tryit}[Labs for this chapter]
Lab~2 (\cref{sec:lab2}) already adds a behavioural right-leg drive to the body model of \cref{ch:interference}.
\end{tryit}

\paragraph{Read with this chapter} Winter and Webster~\cite{winter1983}, Huhta and Webster~\cite{huhta1973}, Webster~\cite{webster}.


\tutpart{Data Converters}{From the quantiser to a measured converter: sampling, specifications and figures of merit, switches, comparators, capacitive DACs, the 12-bit SAR ADC, $\Delta\Sigma$ theory and circuits, converters made for biosignals, and how to test them.}
%% 13_sampling.tex -- Chapter 13: Sampling and Quantisation
%% Status: OUTLINE (build stage 2). Replace the outline with the full chapter when it is written.
\chapter{Sampling and Quantisation}\label{ch:sampling}
\chapterlead{An ADC does two things to a biosignal: it looks at it only at discrete instants, and it rounds every look to one of $2^N$ levels.}

\begin{specnote}[Chapter status: outline]
This chapter is an outline. It will be written in build stage 2 (the Python-only converter core: Chapters 13, 14 and 19 with Labs 6 and 9). The topics, the design question and the labs below are fixed; the derivations, worked examples and figures are still to come.
\end{specnote}

\begin{objectives}
\begin{itemize}
  \item apply the Nyquist theorem to ExG signals and predict what aliases where;
  \item derive the quantisation-noise model, its $6.02N + 1.76$\,dB result and its limits;
  \item use dither and oversampling to trade rate for resolution;
  \item estimate the jitter limit and see why it is minor at biosignal rates.
\end{itemize}
\end{objectives}

\begin{insight}[Design question]
What sample rate and how many bits does each \OB{} mode need, before any circuit is chosen?
\end{insight}

\section{Sampling and aliasing with ExG examples}\label{sec:ch13-1}
Mains harmonics and out-of-band EMG folding into the ECG band; anti-aliasing requirements set by real interferers rather than by the signal.

\begin{keyeq}[Quantisation-limited SNR and effective number of bits]
\begin{equation}
  \mathrm{SQNR} = 6.02N + 1.76\ \mathrm{dB}, \qquad \ENOB = \frac{\mathrm{SNDR} - 1.76}{6.02}
  \label{eq:sqnr}
\end{equation}
\end{keyeq}

\section{The quantisation-noise model}\label{sec:ch13-2}
Uniform error of $\pm\,\mathrm{LSB}/2$, its $\mathrm{LSB}^2/12$ power, and when the white-noise model fails (low-amplitude and DC signals).

\section{Dither and oversampling}\label{sec:ch13-3}
Adding noise to linearise the quantiser; each doubling of the sample rate buys \qty{3}{\deci\bel} of SQNR, setting up the noise shaping of \cref{ch:dstheory}.

\section{Clock jitter}\label{sec:ch13-4}
$\SNR_\mathrm{jitter} = -20\log_{10}(2\pi f_\mathrm{in}\sigma_t)$; at a few kilohertz, even nanoseconds of jitter are harmless, which is why biomedical ADCs worry about energy instead.

\begin{tryit}[Labs for this chapter]
Lab~6 (\cref{sec:lab6}) quantises sine waves, applies coherent sampling and windows, and builds the \code{adctest.py} module every later converter lab uses.
\end{tryit}

\paragraph{Read with this chapter} Pelgrom~\cite{pelgrom}, Schreier, Temes and Pavan~\cite{schreier}.

%% 14_adcspecs.tex -- Chapter 14: Specifications, Figures of Merit and Architecture Choice
%% Status: OUTLINE (build stage 2). Replace the outline with the full chapter when it is written.
\chapter{Specifications, Figures of Merit and Architecture Choice}\label{ch:adcspecs}
\chapterlead{A converter is described by a dozen numbers, but for a biosignal chip two of them, resolution and energy per conversion, decide the architecture.}

\begin{specnote}[Chapter status: outline]
This chapter is an outline. It will be written in build stage 2 (the Python-only converter core: Chapters 13, 14 and 19 with Labs 6 and 9). The topics, the design question and the labs below are fixed; the derivations, worked examples and figures are still to come.
\end{specnote}

\begin{objectives}
\begin{itemize}
  \item define SNR, SNDR, SFDR, THD, ENOB, DNL, INL, offset and gain error and measure each;
  \item compute the Walden and Schreier figures of merit and know when each applies;
  \item place biomedical converters on the Murmann ADC survey;
  \item explain why thermal-noise-limited power grows $\times 4$ per bit and $CV^2$-limited power $\times 2$;
  \item choose between SAR, $\Delta\Sigma$, incremental, VCO-based and event-driven converters for a given biosignal.
\end{itemize}
\end{objectives}

\begin{insight}[Design question]
Which converter architecture fits each \OB{} mode, and what energy per conversion should you expect?
\end{insight}

\section{Dynamic and static specifications}\label{sec:ch14-1}
Spectral metrics from an FFT, and code-density metrics (DNL, INL) from a histogram; offset and gain error and why biomedical systems can often calibrate them away.

\begin{keyeq}[Walden and Schreier figures of merit]
\begin{equation}
  \mathrm{FoM_W} = \frac{P}{f_s\,2^{\ENOB}}, \qquad \mathrm{FoM_S} = \mathrm{SNDR} + 10\log_{10}\!\frac{f_s/2}{P}
  \label{eq:fom}
\end{equation}
\end{keyeq}

\begin{table}[tbp]
\centering\small
\caption{\textbf{Converter architectures for biosignals.} The two converters built in this tutorial are the SAR and the incremental $\Delta\Sigma$; the others appear in \cref{ch:bioadc} or only for contrast.}
\label{tab:ch14-arch}
\begin{tabularx}{\linewidth}{@{}l L L L@{}}
\toprule
Architecture & Typical bio use & Strength & Weakness \\
\midrule
SAR & Multiplexed ExG, neural spikes, 8--12\,b & Lowest energy per conversion, mostly digital, scales with technology & Needs an anti-alias filter and S/H; $kT/C$ and DAC matching limit the bits \\
$\Delta\Sigma$ (DT or CT) & ECG/EEG at 14--20\,b & Relaxed anti-alias filter, high resolution, mismatch-tolerant with a 1-bit DAC & Decimation-filter power; continuous operation; idle tones \\
Incremental $\Delta\Sigma$ & Multiplexed high-resolution channels & $\Delta\Sigma$ accuracy with a reset per sample, multiplexer-friendly & Lower SQNR per OSR than a free-running $\Delta\Sigma$ \\
VCO / time-domain & Direct digitisation, tiny area & Mostly digital, inherent first-order shaping & VCO nonlinearity, supply sensitivity \\
Level-crossing / event-driven & Sparse signals (spikes, ECG) & Data and power scale with activity & Non-uniform samples, harder DSP \\
Pipeline / flash & Not used for ExG & Speed & Power far beyond biosignal needs; shown only for contrast \\
\bottomrule
\end{tabularx}
\end{table}

\section{Figures of merit}\label{sec:ch14-2}
Walden FoM for converters limited by energy per step, Schreier FoM for thermal-noise-limited designs; reading the survey data of Murmann~\cite{murmannsurvey}.

\section{Power per bit}\label{sec:ch14-3}
$kT/C$ sets the capacitor size and therefore the energy for high-resolution converters ($\times 4$ per extra bit); switching energy $CV^2$ dominates at low resolution ($\times 2$ per bit).

\section{Architecture choice for biosignals}\label{sec:ch14-4}
The summary in \cref{tab:ch14-arch}: SAR for multiplexed EMG and spikes, incremental $\Delta\Sigma$ for high-resolution ECG and EEG, and newer time-domain and event-driven options.

\begin{tryit}[Labs for this chapter]
Lab~6 (\cref{sec:lab6}) computes SNDR, SFDR and ENOB and histogram DNL/INL with injected capacitor mismatch.
\end{tryit}

\paragraph{Read with this chapter} Murmann~\cite{murmannsurvey}, Pelgrom~\cite{pelgrom}.

%% 15_switches.tex -- Chapter 15: Sample-and-Hold Circuits and Switches
%% Status: OUTLINE (build stage 4). Replace the outline with the full chapter when it is written.
\chapter{Sample-and-Hold Circuits and Switches}\label{ch:switches}
\chapterlead{The sampling switch and capacitor set the noise floor of a SAR converter before a single comparison is made.}

\begin{specnote}[Chapter status: outline]
This chapter is an outline. It will be written in build stage 4 (the rest of Part III with Labs 7, 8 and 10). The topics, the design question and the labs below are fixed; the derivations, worked examples and figures are still to come.
\end{specnote}

\begin{objectives}
\begin{itemize}
  \item derive $kT/C$ noise and the minimum sampling capacitance per bit;
  \item size a switch for settling within the acquisition time;
  \item reduce charge injection and clock feedthrough with bottom-plate sampling;
  \item design a bootstrapped switch for a \qty{1.2}{\volt} supply;
  \item handle leakage at very low sample rates and multiplexer crosstalk.
\end{itemize}
\end{objectives}

\begin{insight}[Design question]
How large must the sampling capacitor be, and how should the switch be built, for 12 bits at \qty{1.2}{\volt} and a few kilosamples per second?
\end{insight}

\section{\texorpdfstring{$kT/C$}{kT/C} noise and capacitor size}\label{sec:ch15-1}
The sampled thermal noise is independent of the switch resistance; the capacitance needed for $N$ bits grows as $2^{2N}$.

\begin{keyeq}[Sampled noise and minimum capacitance]
\begin{equation}
  v_{n,\mathrm{rms}} = \sqrt{kT/C} \quad\Rightarrow\quad C_{\min} \approx \frac{12\,kT\,2^{2N}}{V_{\mathrm{FS}}^2}
  \label{eq:ktc}
\end{equation}
\end{keyeq}

\section{Settling and on-resistance}\label{sec:ch15-2}
Exponential settling to half an LSB in the acquisition window; the on-resistance across the input range at a low supply.

\section{Charge injection and clock feedthrough}\label{sec:ch15-3}
Signal-dependent errors from channel charge; bottom-plate sampling and dummy switches.

\section{Bootstrapped switches}\label{sec:ch15-4}
Keeping $V_\mathrm{GS}$ constant to linearise the switch at \qty{1.2}{\volt}; reliability limits in \sgx{}.

\section{Leakage and multiplexing}\label{sec:ch15-5}
At a few hundred samples per second the held charge leaks through off switches, a problem specific to biomedical ADCs; crosstalk between multiplexed channels.

\begin{tryit}[Labs for this chapter]
Lab~8 (\cref{sec:lab8}) builds a bootstrapped sample-and-hold as part of the SAR ADC.
\end{tryit}

\paragraph{Read with this chapter} Razavi~\cite{razavi}, Pelgrom~\cite{pelgrom}.

%% 16_comparators.tex -- Chapter 16: Comparators
%% Status: OUTLINE (build stage 4). Replace the outline with the full chapter when it is written.
\chapter{Comparators}\label{ch:comparators}
\chapterlead{A SAR converter makes one comparator decision per bit; its offset, noise and energy set much of the converter's performance.}

\begin{specnote}[Chapter status: outline]
This chapter is an outline. It will be written in build stage 4 (the rest of Part III with Labs 7, 8 and 10). The topics, the design question and the labs below are fixed; the derivations, worked examples and figures are still to come.
\end{specnote}

\begin{objectives}
\begin{itemize}
  \item compare static preamplifier-plus-latch, dynamic StrongARM and double-tail comparators;
  \item predict offset from Pelgrom coefficients and calibrate it;
  \item measure input-referred noise with transient-noise simulation;
  \item control kickback and metastability;
  \item estimate energy per decision.
\end{itemize}
\end{objectives}

\begin{insight}[Design question]
Which comparator gives 12-bit decisions at the lowest energy, and how do you verify its offset and noise?
\end{insight}

\section{Comparator architectures}\label{sec:ch16-1}
Static preamplifier and latch versus the dynamic StrongARM latch and the double-tail comparator for low supplies.

\section{Offset and calibration}\label{sec:ch16-2}
Monte Carlo offset from input-pair mismatch; foreground and background calibration.

\section{Noise}\label{sec:ch16-3}
Input-referred noise of a regenerative latch and how to extract it from the probability of a '1' versus input voltage in transient-noise simulation.

\section{Kickback and metastability}\label{sec:ch16-4}
Charge kicked back onto the DAC; decisions on tiny inputs that take too long, and what they cost an asynchronous SAR.

\section{Energy per decision}\label{sec:ch16-5}
Switching energy of the latch and its clocking; how it scales with supply and speed.

\begin{tryit}[Labs for this chapter]
Lab~7 (\cref{sec:lab7}) designs a StrongARM comparator in \sgx{}: Monte Carlo offset, transient-noise input-referred noise and energy per decision.
\end{tryit}

\paragraph{Read with this chapter} Razavi~\cite{razavi}, the IHP Analog Academy~\cite{ihpacademy} (the dynamic comparator module).

%% 17_cdac.tex -- Chapter 17: Capacitive DACs and References
%% Status: OUTLINE (build stage 4). Replace the outline with the full chapter when it is written.
\chapter{Capacitive DACs and References}\label{ch:cdac}
\chapterlead{The capacitive DAC is both the sample-and-hold and the reference ladder of a SAR converter; its matching sets the linearity.}

\begin{specnote}[Chapter status: outline]
This chapter is an outline. It will be written in build stage 4 (the rest of Part III with Labs 7, 8 and 10). The topics, the design question and the labs below are fixed; the derivations, worked examples and figures are still to come.
\end{specnote}

\begin{objectives}
\begin{itemize}
  \item compare binary-weighted, split (bridge) and C-2C capacitor arrays;
  \item compute the switching energy of the conventional, monotonic, $V_\mathrm{cm}$-based and merged-capacitor schemes;
  \item size the unit capacitor from matching and from $kT/C$;
  \item choose between MOM and MIM capacitors in \sgx{} and lay them out in common centroid;
  \item design reference buffering and decoupling.
\end{itemize}
\end{objectives}

\begin{insight}[Design question]
How small can the unit capacitor be for 12-bit linearity, and which switching scheme minimises the DAC energy?
\end{insight}

\section{Capacitor array topologies}\label{sec:ch17-1}
Binary arrays, split arrays with a bridge capacitor, and C-2C ladders: area, linearity and parasitic sensitivity.

\section{Switching schemes and energy}\label{sec:ch17-2}
Conventional, monotonic, $V_\mathrm{cm}$-based and merged-capacitor switching; the average energy per conversion of each.

\section{Unit-capacitor sizing}\label{sec:ch17-3}
Matching-limited versus $kT/C$-limited size; why matching usually wins for SAR converters above 10 bits.

\section{Capacitors in \texorpdfstring{\sgx{}}{SG13G2}}\label{sec:ch17-4}
MOM and MIM options, their density and matching, and common-centroid layout of the array.

\section{References}\label{sec:ch17-5}
The DAC draws charge from the reference every bit cycle; on-chip buffers versus large decoupling capacitors.

\begin{tryit}[Labs for this chapter]
Lab~8 (\cref{sec:lab8}) builds a CDAC with monotonic switching; Lab~14 (\cref{sec:lab14}) lays it out.
\end{tryit}

\paragraph{Read with this chapter} Pelgrom~\cite{pelgrom}, Moser and Pretl~\cite{jkusar} (an open non-binary 12-bit CDAC).

%% 18_sar.tex -- Chapter 18: The SAR ADC
%% Status: OUTLINE (build stage 4). Replace the outline with the full chapter when it is written.
\chapter{The SAR ADC}\label{ch:sar}
\chapterlead{Binary search in silicon: the 12-bit asynchronous SAR is the converter for the multiplexed EMG and spike modes of \OB{}.}

\begin{specnote}[Chapter status: outline]
This chapter is an outline. It will be written in build stage 4 (the rest of Part III with Labs 7, 8 and 10). The topics, the design question and the labs below are fixed; the derivations, worked examples and figures are still to come.
\end{specnote}

\begin{objectives}
\begin{itemize}
  \item implement the binary search algorithm with a CDAC, a comparator and SAR logic;
  \item choose between synchronous and asynchronous timing;
  \item write the SAR logic in Verilog and synthesise or hand-build it;
  \item add redundancy with non-binary weights and calibrate digitally;
  \item power-gate and duty-cycle the converter between biosignal samples.
\end{itemize}
\end{objectives}

\begin{insight}[Design question]
What SAR design reaches 12 bits at the lowest energy per conversion for eight multiplexed channels?
\end{insight}

\section{The binary search algorithm}\label{sec:ch18-1}
Sampling, $N$ comparison cycles and the DAC update; the timing of a conversion.

\section{Synchronous versus asynchronous timing}\label{sec:ch18-2}
An external fast clock versus self-timed comparisons; why asynchronous SARs dominate low-power designs.

\section{SAR logic}\label{sec:ch18-3}
Verilog for the successive-approximation register, synthesised with the open flow or built from custom cells.

\section{Redundancy and calibration}\label{sec:ch18-4}
Non-binary weights tolerate DAC settling errors and decision errors; digital correction and calibration.

\section{Power gating and duty cycling}\label{sec:ch18-5}
At a few kilosamples per second the converter idles most of the time; leakage then dominates.

\section{Case study: a 12-bit SAR in a seizure-detection SoC}\label{sec:ch18-6}
The \qty{1}{\volt} SAR inside the EEG acquisition SoC of Verma et al.~\cite{verma2010}, read alongside the open SKY130 design of Moser and Pretl~\cite{jkusar}.

\begin{tryit}[Labs for this chapter]
Lab~8 (\cref{sec:lab8}) integrates the S/H, the CDAC and Verilog SAR logic into a 10- to 12-bit SAR and measures its ENOB and FoM with \code{adctest.py}.
\end{tryit}

\paragraph{Read with this chapter} Verma et al.~\cite{verma2010}, Moser and Pretl~\cite{jkusar}, the TWEPP 2024 talk~\cite{twepp2024sar}, the IHP Analog Academy~\cite{ihpacademy}.

%% 19_dstheory.tex -- Chapter 19: $\Delta\Sigma$ Modulation: Theory
%% Status: OUTLINE (build stage 2). Replace the outline with the full chapter when it is written.
\chapter{$\Delta\Sigma$ Modulation: Theory}\label{ch:dstheory}
\chapterlead{Oversample, quantise coarsely and push the quantisation noise out of the band: noise shaping turns a 1-bit quantiser into a 16-bit converter.}

\begin{specnote}[Chapter status: outline]
This chapter is an outline. It will be written in build stage 2 (the Python-only converter core: Chapters 13, 14 and 19 with Labs 6 and 9). The topics, the design question and the labs below are fixed; the derivations, worked examples and figures are still to come.
\end{specnote}

\begin{objectives}
\begin{itemize}
  \item derive the noise transfer function (NTF) and signal transfer function (STF) of first- and second-order loops;
  \item predict SQNR from loop order and oversampling ratio (OSR);
  \item check stability with the Lee criterion and understand idle tones and dither;
  \item synthesise and realise an NTF with \code{synthesizeNTF} and \code{realizeNTF};
  \item know when linearised analysis fails and simulation is needed.
\end{itemize}
\end{objectives}

\begin{insight}[Design question]
Which loop order and OSR give 16 bits in the ECG/EEG band at the lowest clock rate?
\end{insight}

\section{Noise shaping}\label{sec:ch19-1}
The linear model of a $\Delta\Sigma$ loop; $\mathrm{NTF}(z) = (1-z^{-1})^L$ and its high-pass shape.

\begin{keyeq}[Peak SQNR of an order-$L$ modulator]
\begin{equation}
  \mathrm{SQNR}_{L} \approx 6.02N + 1.76 - 10\log_{10}\!\frac{\pi^{2L}}{2L+1} + (20L+10)\log_{10}\mathrm{OSR}\ \mathrm{dB}
  \label{eq:dssqnr}
\end{equation}
\end{keyeq}

\section{SQNR versus order and OSR}\label{sec:ch19-2}
Each doubling of OSR buys $6L + 3$\,dB; diminishing returns and the trade-off with stability.

\section{Stability, idle tones and dither}\label{sec:ch19-3}
The Lee criterion $\lVert \mathrm{NTF}\rVert_\infty < 1.5$ for 1-bit loops; tones with DC inputs, which matter for ECG baselines.

\section{NTF synthesis}\label{sec:ch19-4}
Using python-deltasigma~\cite{pydeltasigma}, the Python port of Schreier's toolbox~\cite{schreier}, to place zeros and poles and map them to a loop topology.

\begin{tryit}[Labs for this chapter]
Lab~9 (\cref{sec:lab9}) synthesises first- to third-order NTFs, simulates SQNR versus OSR and the incremental mode, and designs a fixed-point CIC decimator.
\end{tryit}

\paragraph{Read with this chapter} Schreier, Temes and Pavan~\cite{schreier}, python-deltasigma~\cite{pydeltasigma}.

%% 20_dscircuits.tex -- Chapter 20: $\Delta\Sigma$ Circuits: Discrete-Time, Continuous-Time and Incremental
%% Status: OUTLINE (build stage 4). Replace the outline with the full chapter when it is written.
\chapter{$\Delta\Sigma$ Circuits: Discrete-Time, Continuous-Time and Incremental}\label{ch:dscircuits}
\chapterlead{From the loop filter on paper to integrators in silicon, and to the incremental converter that serves the high-resolution \OB{} channels.}

\begin{specnote}[Chapter status: outline]
This chapter is an outline. It will be written in build stage 4 (the rest of Part III with Labs 7, 8 and 10). The topics, the design question and the labs below are fixed; the derivations, worked examples and figures are still to come.
\end{specnote}

\begin{objectives}
\begin{itemize}
  \item build switched-capacitor integrators and set their OTA requirements;
  \item design continuous-time loops with resistive or Gm-C integrators and use their inherent anti-aliasing;
  \item handle excess loop delay;
  \item design an incremental $\Delta\Sigma$ converter for multiplexed channels;
  \item design and budget the decimation filter (CIC/sinc$^K$ and half-band FIR).
\end{itemize}
\end{objectives}

\begin{insight}[Design question]
How do you implement the 2nd-order incremental modulator of \OB{} and its decimator within the power budget?
\end{insight}

\section{Switched-capacitor integrators}\label{sec:ch20-1}
Charge transfer, OTA gain and slew requirements, and $kT/C$ noise of the sampling capacitor.

\section{Continuous-time \texorpdfstring{$\Delta\Sigma$}{Delta-Sigma}}\label{sec:ch20-2}
Resistive and Gm-C integrators; the implicit anti-aliasing of a CT loop; clock-jitter sensitivity.

\section{Excess loop delay}\label{sec:ch20-3}
How a fraction of a clock period of delay destabilises the loop and how to compensate it.

\section{Incremental \texorpdfstring{$\Delta\Sigma$}{Delta-Sigma}}\label{sec:ch20-4}
Resetting the loop for every sample so channels can be multiplexed; the resolution versus number of cycles.

\section{Decimation}\label{sec:ch20-5}
CIC (sinc$^K$) and half-band FIR stages, their word lengths and their power.

\section{The reference design}\label{sec:ch20-6}
The full 2nd-order incremental modulator of the \OB{} ECG/EEG mode.

\begin{tryit}[Labs for this chapter]
Lab~10 (\cref{sec:lab10}) implements the Lab~9 loop with switched-capacitor integrators and a 1-bit quantiser in \sgx{} and compares it with the ideal model.
\end{tryit}

\paragraph{Read with this chapter} Schreier, Temes and Pavan~\cite{schreier}.

%% 21_bioadc.tex -- Chapter 21: Biomedical Converter Innovations
%% Status: OUTLINE (build stage 4). Replace the outline with the full chapter when it is written.
\chapter{Biomedical Converter Innovations}\label{ch:bioadc}
\chapterlead{In recent designs the boundary between the amplifier and the converter is dissolving.}

\begin{specnote}[Chapter status: outline]
This chapter is an outline. It will be written in build stage 4 (the rest of Part III with Labs 7, 8 and 10). The topics, the design question and the labs below are fixed; the derivations, worked examples and figures are still to come.
\end{specnote}

\begin{objectives}
\begin{itemize}
  \item explain direct digitisation with $\Delta\Sigma$ front-ends that take the electrode as input;
  \item design for the large input range that motion artefacts require;
  \item describe VCO- and time-based ADCs;
  \item compare level-crossing and event-driven sampling with uniform sampling;
  \item assess compressive sensing for biosignals.
\end{itemize}
\end{objectives}

\begin{insight}[Design question]
When does it pay to merge the amplifier and the ADC, and when to sample only on events?
\end{insight}

\section{Direct digitisation}\label{sec:ch21-1}
$\Delta\Sigma$ loops that close around the electrode, removing the LNA and gaining dynamic range against artefacts.

\section{VCO- and time-domain ADCs}\label{sec:ch21-2}
Voltage-controlled oscillators as integrating quantisers with first-order shaping; nonlinearity and supply sensitivity.

\section{Level-crossing and event-driven sampling}\label{sec:ch21-3}
Data and power that scale with activity for sparse spikes and ECG.

\section{Compressive sensing}\label{sec:ch21-4}
Random projections of sparse signals; what it saves in the radio and costs in reconstruction.

\section{The dissolving AFE/ADC boundary}\label{sec:ch21-5}
Reading recent ISSCC designs with the budgets of this tutorial.

\paragraph{Read with this chapter} Murmann~\cite{murmannsurvey}.

%% 22_adctest.tex -- Chapter 22: Testing and Characterising ADCs
%% Status: OUTLINE (build stage 4). Replace the outline with the full chapter when it is written.
\chapter{Testing and Characterising ADCs}\label{ch:adctest}
\chapterlead{A converter is only as good as the measurement that proves it.}

\begin{specnote}[Chapter status: outline]
This chapter is an outline. It will be written in build stage 4 (the rest of Part III with Labs 7, 8 and 10). The topics, the design question and the labs below are fixed; the derivations, worked examples and figures are still to come.
\end{specnote}

\begin{objectives}
\begin{itemize}
  \item choose coherent sampling frequencies and windows for FFT tests;
  \item compute SNDR and SFDR from an FFT;
  \item measure DNL and INL with the sine-wave histogram and with static ramps;
  \item compare behavioural Verilog-A or Python models with transistor-level results;
  \item plan the measurements for a returned OpenMPW die.
\end{itemize}
\end{objectives}

\begin{insight}[Design question]
How do you prove, in simulation and on a returned die, that the \OB{} converters meet their specifications?
\end{insight}

\section{Coherent sampling and windowing}\label{sec:ch22-1}
Choosing $f_\mathrm{in} = (M/N)\fs$ with $M$ and $N$ co-prime; windows when coherence is impossible.

\section{FFT-based SNDR and SFDR}\label{sec:ch22-2}
Bins, leakage, harmonics and the noise floor; the processing gain of a long FFT.

\section{Histogram and ramp tests}\label{sec:ch22-3}
Code-density DNL and INL from a sine-wave histogram; static ramp tests for slow converters.

\section{Models versus transistors}\label{sec:ch22-4}
Behavioural models for fast system simulation and their validation against transistor-level results.

\section{Testing a returned die}\label{sec:ch22-5}
The bench set-up, the measurements that matter, and how to report them.

\begin{tryit}[Labs for this chapter]
Lab~6 (\cref{sec:lab6}) writes the \code{adctest.py} module used by Labs~8 and~10.
\end{tryit}

\paragraph{Read with this chapter} Pelgrom~\cite{pelgrom}.


\tutpart{Biosignal Processing and Systems}{What happens to the samples: a DSP toolkit for biosignals, ECG, EEG, EMG and spike processing on open datasets, on-chip feature extraction and system blocks, and bioimpedance sensing with IQ demodulation.}
%% 23_dsp.tex -- Chapter 23: A DSP Toolkit for Biosignals
%% Status: OUTLINE (build stage 5). Replace the outline with the full chapter when it is written.
\chapter{A DSP Toolkit for Biosignals}\label{ch:dsp}
\chapterlead{Most of what a clinician sees has been filtered, decimated and cleaned after the ADC. The filters must not distort the waveform features they are meant to reveal.}

\begin{specnote}[Chapter status: outline]
This chapter is an outline. It will be written in build stage 5 (Part IV, including bioimpedance, with Labs 11--13 and 15). The topics, the design question and the labs below are fixed; the derivations, worked examples and figures are still to come.
\end{specnote}

\begin{objectives}
\begin{itemize}
  \item choose sample rates for each biosignal;
  \item compare FIR and IIR filters and use zero-phase filtering;
  \item remove baseline wander and mains interference with adaptive and comb notch filters;
  \item use wavelets for denoising and feature extraction;
  \item predict fixed-point effects and choose word lengths.
\end{itemize}
\end{objectives}

\begin{insight}[Design question]
Which filters, at what word length, clean each biosignal without distorting it, and how much do they cost on chip?
\end{insight}

\section{Sample-rate choices}\label{sec:ch23-1}
Rates for diagnostic and monitoring ECG, EEG, EMG and spikes, and decimation between them.

\section{FIR, IIR and zero-phase filtering}\label{sec:ch23-2}
Phase distortion of ST segments by causal IIR high-pass filters; forward--backward filtering offline.

\section{Baseline wander and mains interference}\label{sec:ch23-3}
High-pass filtering versus spline and median baselines; adaptive notch and comb filters that track \qty{50}{\hertz} or \qty{60}{\hertz} and their harmonics.

\section{Wavelets}\label{sec:ch23-4}
Multi-resolution decompositions for denoising and QRS detection.

\section{Fixed-point effects}\label{sec:ch23-5}
Coefficient quantisation, limit cycles and word-length choices for on-chip filters.

\paragraph{Read with this chapter} Tompkins~\cite{tompkins}.

%% 24_ecg.tex -- Chapter 24: ECG Processing
%% Status: OUTLINE (build stage 5). Replace the outline with the full chapter when it is written.
\chapter{ECG Processing}\label{ch:ecg}
\chapterlead{Detecting heartbeats is the oldest biomedical DSP problem, and the best way to see how ADC bits and amplifier noise reach a clinical result.}

\begin{specnote}[Chapter status: outline]
This chapter is an outline. It will be written in build stage 5 (Part IV, including bioimpedance, with Labs 11--13 and 15). The topics, the design question and the labs below are fixed; the derivations, worked examples and figures are still to come.
\end{specnote}

\begin{objectives}
\begin{itemize}
  \item implement the Pan--Tompkins QRS detector;
  \item compute heart rate and heart-rate variability metrics;
  \item classify beats on the MIT-BIH Arrhythmia Database;
  \item test robustness with the MIT-BIH Noise Stress Test Database;
  \item show how ADC bits and AFE noise move detection sensitivity.
\end{itemize}
\end{objectives}

\begin{insight}[Design question]
How many bits and how much input noise can the ECG chain afford before QRS detection degrades?
\end{insight}

\section{Pan--Tompkins QRS detection}\label{sec:ch24-1}
Band-pass, derivative, squaring, moving-window integration and adaptive thresholds~\cite{pantompkins1985}.

\section{Heart rate and HRV}\label{sec:ch24-2}
RR intervals, time-domain and frequency-domain HRV metrics.

\section{Beat classification}\label{sec:ch24-3}
Features and simple classifiers on the MIT-BIH Arrhythmia Database~\cite{mitbih}.

\section{Noise stress testing}\label{sec:ch24-4}
Calibrated baseline wander, muscle and electrode-motion noise from the Noise Stress Test Database~\cite{nstdb}.

\section{Co-design: bits and noise versus sensitivity}\label{sec:ch24-5}
Sweeping the front-end parameters of a behavioural chain against detection sensitivity and positive predictivity.

\begin{tryit}[Labs for this chapter]
Lab~11 (\cref{sec:lab11}) passes MIT-BIH and Noise Stress Test records through a behavioural AFE and ADC and plots QRS sensitivity and PPV versus ADC bits and input noise.
\end{tryit}

\paragraph{Read with this chapter} Pan and Tompkins~\cite{pantompkins1985}, Tompkins~\cite{tompkins}, Moody and Mark~\cite{mitbih}.

%% 25_eeg.tex -- Chapter 25: EEG Processing and Brain--Computer Interfaces
%% Status: OUTLINE (build stage 5). Replace the outline with the full chapter when it is written.
\chapter{EEG Processing and Brain--Computer Interfaces}\label{ch:eeg}
\chapterlead{EEG carries information in rhythms and their changes; most of the processing estimates power in frequency bands.}

\begin{specnote}[Chapter status: outline]
This chapter is an outline. It will be written in build stage 5 (Part IV, including bioimpedance, with Labs 11--13 and 15). The topics, the design question and the labs below are fixed; the derivations, worked examples and figures are still to come.
\end{specnote}

\begin{objectives}
\begin{itemize}
  \item compute band power and spectral features;
  \item remove EOG and EMG artefacts, including with ICA;
  \item detect seizure onset on the CHB-MIT database;
  \item explain the basics of SSVEP and motor-imagery BCIs.
\end{itemize}
\end{objectives}

\begin{insight}[Design question]
Which EEG features need which resolution and sample rate, and what does that save on chip?
\end{insight}

\section{Band power and spectral features}\label{sec:ch25-1}
Welch spectra, relative band power and spectral edge frequency for the $\delta$ to $\gamma$ bands of \cref{ch:signals}.

\section{Artefact removal}\label{sec:ch25-2}
Eye-blink and muscle artefacts; regression and independent component analysis.

\section{Seizure-onset detection}\label{sec:ch25-3}
Patient-specific detectors on the CHB-MIT database~\cite{chbmit} and the SoC of Verma et al.~\cite{verma2010}.

\section{BCI basics}\label{sec:ch25-4}
Steady-state visually evoked potentials and motor imagery.

\begin{tryit}[Labs for this chapter]
Lab~12 (\cref{sec:lab12}) computes band-power features on CHB-MIT, trains a simple classifier and repeats it with reduced word length and sample rate.
\end{tryit}

\paragraph{Read with this chapter} Shoeb~\cite{chbmit}, Verma et al.~\cite{verma2010}.

%% 26_emg.tex -- Chapter 26: EMG and Neural Spike Processing
%% Status: OUTLINE (build stage 5). Replace the outline with the full chapter when it is written.
\chapter{EMG and Neural Spike Processing}\label{ch:emg}
\chapterlead{Muscle signals are read through their envelope; neural signals through the timing of individual spikes.}

\begin{specnote}[Chapter status: outline]
This chapter is an outline. It will be written in build stage 5 (Part IV, including bioimpedance, with Labs 11--13 and 15). The topics, the design question and the labs below are fixed; the derivations, worked examples and figures are still to come.
\end{specnote}

\begin{objectives}
\begin{itemize}
  \item compute EMG envelopes and the time-domain features MAV, WL, ZC and SSC;
  \item classify hand gestures on NinaPro;
  \item detect neural spikes with the nonlinear energy operator (NEO);
  \item explain spike sorting and why implants compress before transmitting.
\end{itemize}
\end{objectives}

\begin{insight}[Design question]
How much resolution, sample rate and channel count does gesture recognition or spike detection actually need?
\end{insight}

\section{EMG envelopes and features}\label{sec:ch26-1}
Rectification and smoothing; mean absolute value, waveform length, zero crossings and slope-sign changes.

\section{Gesture classification}\label{sec:ch26-2}
LDA and SVM classifiers on the NinaPro database~\cite{ninapro}.

\section{Spike detection}\label{sec:ch26-3}
Thresholds and the NEO; detection versus the noise of \cref{ch:noise}.

\section{Spike sorting and compression}\label{sec:ch26-4}
Clustering waveforms; why a thousand-channel implant must send spikes, not samples.

\begin{tryit}[Labs for this chapter]
Lab~13 (\cref{sec:lab13}) computes MAV, WL, ZC and SSC on NinaPro, classifies with LDA or SVM and repeats with fixed-point features.
\end{tryit}

\paragraph{Read with this chapter} Atzori et al.~\cite{ninapro}.

%% 27_soc.tex -- Chapter 27: On-Chip Processing and System Integration
%% Status: OUTLINE (build stage 5). Replace the outline with the full chapter when it is written.
\chapter{On-Chip Processing and System Integration}\label{ch:soc}
\chapterlead{On a wearable or implant, the radio usually costs more energy than the sensing; computing on chip is how you spend less of it.}

\begin{specnote}[Chapter status: outline]
This chapter is an outline. It will be written in build stage 5 (Part IV, including bioimpedance, with Labs 11--13 and 15). The topics, the design question and the labs below are fixed; the derivations, worked examples and figures are still to come.
\end{specnote}

\begin{objectives}
\begin{itemize}
  \item place feature extraction and small machine-learning models in silicon;
  \item use the seizure-detection SoC of Verma et al.~\cite{verma2010} as a case study;
  \item compress biosignal data before transmission;
  \item design power management with LDOs and energy harvesting;
  \item compare BLE, backscatter and inductive data links;
  \item introduce stimulation: current DACs and charge balance.
\end{itemize}
\end{objectives}

\begin{insight}[Design question]
Which processing belongs on the \OB{} chip, and which data must leave it?
\end{insight}

\section{Feature extraction and edge ML in silicon}\label{sec:ch27-1}
Energy per operation versus energy per transmitted bit.

\section{Case study: an EEG seizure-detection SoC}\label{sec:ch27-2}
Chopper IA, 12-bit SAR and feature processor on one die~\cite{verma2010}.

\section{Compression}\label{sec:ch27-3}
Lossless and lossy compression for ECG, EEG and spikes.

\section{Power management}\label{sec:ch27-4}
LDOs, their noise and PSRR for microvolt front-ends; batteries and energy harvesting.

\section{Data links}\label{sec:ch27-5}
BLE, backscatter and inductive links; energy per bit.

\section{Stimulation}\label{sec:ch27-6}
Current-steering DACs, charge balance and electrode safety.

\begin{tryit}[Labs for this chapter]
Lab~12 (\cref{sec:lab12}) relates seizure-detection accuracy to bits and a power proxy.
\end{tryit}

\paragraph{Read with this chapter} Verma et al.~\cite{verma2010}, Sarpeshkar~\cite{sarpeshkar}.

%% 28_bioz.tex -- Chapter 28: Bioimpedance Sensing and IQ Demodulation
%% Status: OUTLINE (build stage 5). Replace the outline with the full chapter when it is written.
\chapter{Bioimpedance Sensing and IQ Demodulation}\label{ch:bioz}
\chapterlead{Inject a small AC current, measure the voltage, and recover the complex impedance of tissue or of the electrode interface.}

\begin{specnote}[Chapter status: outline]
This chapter is an outline. It will be written in build stage 5 (Part IV, including bioimpedance, with Labs 11--13 and 15). The topics, the design question and the labs below are fixed; the derivations, worked examples and figures are still to come.
\end{specnote}

\begin{objectives}
\begin{itemize}
  \item model tissue and electrodes with the Cole model and its dispersion;
  \item explain why a tetrapolar (four-electrode) measurement removes electrode impedance;
  \item design excitation sources: sine DAC with V-to-I converter, Howland and OTA-based current sources, square and pseudo-sine waves, multi-sine and chirp;
  \item implement IQ (quadrature) demodulation in the analog and digital domains and predict its errors;
  \item run ECG and impedance on the same electrodes by frequency separation;
  \item respect the frequency-dependent patient-current limits.
\end{itemize}
\end{objectives}

\begin{insight}[Design question]
What excitation frequency, current and demodulator give the required $\Delta Z$ resolution within the safety limits?
\end{insight}

\section{Applications}\label{sec:ch28-1}
Respiration by impedance pneumography (reusing the ECG electrodes), impedance cardiography, body composition (BIA/BIS), electrical impedance tomography (EIT), electrodermal activity, electrode-contact quality and cell and tissue impedance spectroscopy.

\begin{keyeq}[IQ demodulation]
\begin{equation}
  I = \tfrac{A}{2}\cos\varphi,\quad Q = \tfrac{A}{2}\sin\varphi \quad\Rightarrow\quad |Z| = \frac{2\sqrt{I^2+Q^2}}{I_{\mathrm{exc}}},\quad \varphi = \operatorname{atan2}(Q, I)
  \label{eq:iq}
\end{equation}
\end{keyeq}

\section{Tissue and electrode models}\label{sec:ch28-2}
The Cole model and dispersion; bipolar versus tetrapolar measurement; the small $\Delta Z$ of respiration and heartbeat on top of a large baseline $Z$, which sets the dynamic range.

\section{Excitation}\label{sec:ch28-3}
Sine DAC plus V-to-I converter, Howland and OTA-based current sources, square-wave and pseudo-sine excitation, multi-sine and chirp for spectroscopy; output impedance, common-mode control and the frequency-dependent patient-current limits of \cref{ch:interference}.

\section{IQ demodulation}\label{sec:ch28-4}
Mixing with in-phase and quadrature references, low-pass filtering, magnitude and phase; analog mixers versus chopper-style square-wave demodulators (and their odd-harmonic errors) versus digital IQ after the ADC; reference phase errors and I/Q mismatch; synchronous detection as a lock-in amplifier.

\section{Readout architecture}\label{sec:ch28-5}
IA and baseline cancellation, ADC requirements for the IQ outputs, ECG and impedance on the same electrodes, and multi-frequency, multi-channel systems for EIT.

\begin{tryit}[Labs for this chapter]
Lab~15 (\cref{sec:lab15}) models Cole tissue in a tetrapolar set-up with digital IQ demodulation in Python, then builds an OTA current source and a square-wave IQ demodulator in \sgx{}.
\end{tryit}

\paragraph{Read with this chapter} Webster~\cite{webster}, Pall\`as-Areny and Webster~\cite{pallasareny}.


\tutpart{Putting It Together}{The full Open-BioAFE architecture and its budgets, the open-source design and verification flow from schematic to an IHP OpenMPW submission, and fifteen hands-on labs with a capstone.}
%% 29_arch.tex -- Chapter 29: The Full Architecture and Its Budgets
%% Status: OUTLINE (build stage 6). Replace the outline with the full chapter when it is written.
\chapter{The Full Architecture and Its Budgets}\label{ch:arch}
\chapterlead{Every chapter so far has answered one budget question. This chapter puts the answers together on one chip.}

\begin{specnote}[Chapter status: outline]
This chapter is an outline. It will be written in build stage 6 (Part V, Lab 14, the capstone and the appendices). The topics, the design question and the labs below are fixed; the derivations, worked examples and figures are still to come.
\end{specnote}

\begin{objectives}
\begin{itemize}
  \item assemble the noise, CMRR, power, area and data-rate budgets of \OB{};
  \item choose SAR or $\Delta\Sigma$ per mode;
  \item add the optional bioimpedance channel;
  \item draw the top-level floorplan.
\end{itemize}
\end{objectives}

\begin{insight}[Design question]
Does the assembled \OB{} chip meet every number of its specification (\cref{tab:ch1-spec}), and where is the margin?
\end{insight}

\section{Noise budget}\label{sec:ch29-1}
From electrode noise (\cref{ch:electrodes}) through the LNA, PGA and ADC to the \qty{1.5}{\uVrms} EEG target.

\section{CMRR budget}\label{sec:ch29-2}
Electrode imbalance, amplifier mismatch and right-leg drive against the \qty{90}{\deci\bel} target.

\section{Power, area and data rate}\label{sec:ch29-3}
The \qty{10}{\micro\watt} per channel front-end budget, converter energy, and the output data rate per mode.

\section{SAR or \texorpdfstring{$\Delta\Sigma$}{Delta-Sigma} per mode}\label{sec:ch29-4}
The routing of \cref{fig:ch1-chain} justified by the budgets.

\section{The bioimpedance channel}\label{sec:ch29-5}
The optional respiration channel of \cref{ch:bioz}.

\section{Floorplan}\label{sec:ch29-6}
Placing eight channels, two converters, references and the digital back end.

\paragraph{Read with this chapter} Verma et al.~\cite{verma2010}, Yazicioglu, Van Hoof and Puers~\cite{yazicioglubook}.

%% 30_flow.tex -- Chapter 30: Open-Source Design and Verification Flow
%% Status: OUTLINE (build stage 6). Replace the outline with the full chapter when it is written.
\chapter{Open-Source Design and Verification Flow}\label{ch:flow}
\chapterlead{Every circuit in this tutorial can be designed, verified and taped out with free tools and an open PDK.}

\begin{specnote}[Chapter status: outline]
This chapter is an outline. It will be written in build stage 6 (Part V, Lab 14, the capstone and the appendices). The topics, the design question and the labs below are fixed; the derivations, worked examples and figures are still to come.
\end{specnote}

\begin{objectives}
\begin{itemize}
  \item set up IIC-OSIC-TOOLS with Xschem, ngspice/Xyce and KLayout;
  \item write and simulate SAR logic in Verilog with mixed-signal co-simulation;
  \item run corners and Monte Carlo;
  \item run DRC, LVS and parasitic extraction in KLayout;
  \item use the JKU AMS chip template and submit to an IHP OpenMPW shuttle.
\end{itemize}
\end{objectives}

\begin{insight}[Design question]
How do you take an \OB{} block from schematic to a verified GDS that IHP will accept?
\end{insight}

\section{The tool set}\label{sec:ch30-1}
IIC-OSIC-TOOLS~\cite{iicosic}: Xschem, ngspice, Xyce, KLayout, Verilator and Icarus Verilog in one container, with the IHP-Open-PDK~\cite{ihppdk}.

\section{Schematic and simulation}\label{sec:ch30-2}
Xschem testbenches; ngspice and Xyce for AC, noise, transient and transient-noise analyses.

\section{Mixed-signal co-simulation}\label{sec:ch30-3}
Verilog SAR logic with the analog CDAC and comparator.

\section{Corners and Monte Carlo}\label{sec:ch30-4}
Process, voltage and temperature corners; mismatch; what a biomedical chip must survive.

\section{Layout and verification}\label{sec:ch30-5}
KLayout DRC, LVS and PEX; post-layout simulation.

\section{From template to tape-out}\label{sec:ch30-6}
The JKU ihp-sg13g2-ams-chip-template~\cite{jkutemplate} with its padframe, and the IHP OpenMPW route documented in IHP-Open-DesignLib~\cite{ihpdesignlib}.

\begin{tryit}[Labs for this chapter]
Lab~14 (\cref{sec:lab14}) lays out the comparator and CDAC, runs DRC, LVS and PEX with the JKU template and compares post-layout with schematic simulation.
\end{tryit}

\paragraph{Read with this chapter} Dorrer and Pretl~\cite{jkutemplate}, IIC-OSIC-TOOLS~\cite{iicosic}, IHP-Open-DesignLib~\cite{ihpdesignlib}.

\include{tutorial/chapters/31_labs}

\appendix
\include{tutorial/chapters/A_appendices}

\backmatter
\cleardoublepage
\renewcommand{\bibname}{References}
\phantomsection\addcontentsline{toc}{chapter}{References}
\bibliographystyle{unsrtnat}
\bibliography{tutorial/bib/core}
\end{document}
